\documentclass[11pt]{article}
\usepackage{preamble}
\setlength{\parskip}{4pt}

\begin{document}

\daybanner{AP Statistics \textperiodcentered\ Unit 3 \textperiodcentered\ Topic 3.14 \textperiodcentered\ B: 2026-12-21 \textperiodcentered\ E: 2027-02-03 \textperiodcentered\ TEACHER KEY}{A Two-Way Table Before Inference}

\sectionbanner{CED TETHER}

{\small
\begin{itemize}[leftmargin=1.5em,itemsep=2pt,topsep=2pt]
  \item \textbf{Skill 3.A} --- Determine relative frequencies, proportions, or probabilities using simulation or calculations.
  \item \textbf{EU VAR-8} --- The chi-square distribution may be used to model variation.
  \item \textbf{LO VAR-8.H} --- Calculate expected counts for two-way tables of categorical data. [Skill 3.A]
  \item \textbf{EK VAR-8.H.1} --- The expected count in a particular cell of a two-way table of categorical data can be calculated using the formula: expected count = (row total)(column total) / table total.
\end{itemize}
}
\textbf{Essential question:} Why is a chi-square condition judged from expected counts rather than the observed cells?\par
\textbf{DOK-3 skill:} 4.C \textperiodcentered\ \textbf{FRQ pattern:} {\small\texttt{expected-\allowbreak{}count-\allowbreak{}boundary-\allowbreak{}and-\allowbreak{}remedy}}\par
\textbf{DOK rationale:} Part (c) coordinates the null-table calculation, a strict boundary condition, sampling evidence, and the meaning of possible remedies to adjudicate two claims and explain why the tempting observed-count rule is unreliable.\par

\frameworkphaseheader{Do Now (first take) $\rightarrow$ Explore (video + follow-along) $\rightarrow$ Exit (finish + turn in)}{1 $\rightarrow$ 3}{45}{%
\begin{itemize}[leftmargin=1.5em,itemsep=2pt,topsep=2pt]
  \item Board slide up at the bell. Ask which evidence is needed before deciding whether the approximation is ready.
  \item Begin the 8.4 video follow-along at minute 5.
  \item Return at minute 33. Require the margins that generate the boundary cell and a defensible next step.
  \item Collect at the door. Score part (c) only, E/P/I.
\end{itemize}
}{%
\begin{itemize}[leftmargin=1.5em,itemsep=2pt,topsep=2pt]
  \item Commit to one claim and identify the evidence still needed.
  \item Complete the follow-along, finish (a)--(c), revise, and turn in.
\end{itemize}
}{%
\begin{itemize}[leftmargin=1.5em,itemsep=2pt,topsep=2pt]
  \item Which row total and column total generate the Not done/Phone cell?
  \item Why can an expected count equal 5 even when its observed count is 6?
  \item When would combining support channels preserve the research question?
\end{itemize}
}{Solo teacher. Require the expected table, the strict boundary check, and a remedy that preserves the question.}

\begin{callout}{\IconWarn\ WATCH FOR}{calloutred}
\begin{itemize}[leftmargin=1.5em,itemsep=2pt,topsep=2pt]
  \item Using the observed interior count as its own expected count.
  \item Rounding an expected count of exactly 5 upward before checking the condition.
  \item Combining categories solely to make the condition pass without considering meaning.
\end{itemize}
\end{callout}

\sectionbanner{SCORING PART (c) --- E / P / I}

\textbf{Expected elements}
\begin{itemize}[leftmargin=1.5em,itemsep=2pt,topsep=2pt]
  \item \textbf{selects-expected-count-rule} (required) --- Uses expected, not observed, counts
  \item \textbf{uses-boundary-values} (required) --- Contrasts observed 6 with expected 5 and applies $>5$ strictly
  \item \textbf{withholds-usual-approximation} (required) --- Does not proceed under the stated rule
  \item \textbf{states-reader-consequence} (required) --- Explains why observed counts cannot validate the null approximation
  \item \textbf{gives-honest-remedy} (required) --- Gives a remedy that preserves the research question
\end{itemize}

{\small\renewcommand{\arraystretch}{1.25}
\noindent\begin{tabular}{|p{0.5in}|p{5.9in}|}
\hline
\textbf{E} & Uses expected 5 to reject the usual approximation, diagnoses the observed-count rule, and gives an honest remedy. \\ \hline
\textbf{P} & Finds the boundary failure but omits the consequence, strict comparison, or defensible remedy. \\ \hline
\textbf{I} & Uses observed counts, treats 5 as greater than 5, or merges categories only to pass. \\ \hline
\end{tabular}}

\clearpage
\focusbanner{TODAY'S PROBLEM --- annotated}

Suppose a membership association has 5,000 new members and takes an SRS of 120. Each member is classified by orientation status and preferred support channel. The association may test whether the variables are associated.\par\smallskip \textbf{Nico:} ``Every observed cell is at least 6, and the data came from an SRS, so the usual chi-square setup is ready.''\par \textbf{Priya:} ``The margins may produce a different boundary check; I would calculate that table before deciding.''\par
\begin{center}
\textbf{Support channel}\par\smallskip
\begin{tabular}{|l|c|c|c|c|}
\hline
\textbf{Status} & \textbf{Chat} & \textbf{Email} & \textbf{Phone} & \textbf{Total} \\ \hline
\textbf{Done} & 42 & 34 & 24 & 100 \\ \hline
\textbf{Not done} & 8 & 6 & 6 & 20 \\ \hline
\textbf{Total} & 50 & 40 & 30 & 120 \\ \hline
\end{tabular}
\end{center}
\begin{firsttakebox}{FIRST TAKE (5 min)}
Which claim is more defensible so far? Identify the evidence that should settle the condition check.\par
\answer{Not graded. The observed cells were chosen to pass a tempting but incorrect checklist while one expected cell lands exactly on the boundary.}
\end{firsttakebox}

\textit{Rules callout ``EXPECTED COUNTS UNDER NO ASSOCIATION (keep on the page)'' is printed on the student sheet.}\par\medskip

\dokbadge{1}~\textbf{(a)} Use the margins to calculate the expected counts for Done/Chat and Not done/Phone under no association.\par
\answer{Done/Chat: $E=(100)(50)/120=41.6667$.\par Not done/Phone: $E=(20)(30)/120=5.0000$.}

\dokbadge{2}~\textbf{(b)} Calculate all six expected counts, identify the smallest, and state the degrees of freedom. Verify the random-sample and 10 percent conditions.\par
\answer{Expected counts are Done: $(41.6667,33.3333,25.0000)$;\par Not done: $(8.3333,6.6667,5.0000)$. The minimum is 5 and $df=2$. The SRS condition holds, and $120\leq0.10(5{,}000)=500$.}

\dokbadge{3}~\textbf{(c)} Adjudicate the claims and decide whether the usual chi-square approximation should proceed as is. Cite the observed and expected boundary counts, explain what the rejected rule would mislead a reader into believing, and recommend a defensible next step without silently changing the research question.\par
\answer{Priya is warranted. Although the observed Not-done/Phone count is 6, its expected count is $(20)(30)/120=5$, not greater than 5. Under the stated conservative rule, the usual approximation should not proceed as is. The rejected observed-count rule would falsely suggest that the realized table establishes the null approximation's accuracy. Collect a larger SRS and recheck every expected count. Combine channels only for a substantive, preferably pre-data reason; otherwise the remedy changes the research question and may hide a pattern.}

\begin{summaryexitbox}{EXIT}
One thing the video changed about my first take: \hrulefill
\end{summaryexitbox}

\end{document}
